\subsection{An exponent-thirteen Prym and its central value}
\label{subsec:thirteen}

The complementary rank-zero target at exponent thirteen is already
implicit in \citet[Proposition~2.1 and Corollary~2.2, p.~942]{Tzermias2004},
using the rank-one calculation of Gross--Rohrlich. The decomposition
over the cyclotomic field is due to
\citet[Theorem~6.2 and Corollary~6.3, pp.~89--90]{Lim1992}.
We give explicit quotient equations, the resulting Hecke constituents,
and a central-value certificate for this classical target. Its Selmer
input is proved separately in Section~\ref{subsec:lm-lz-rational}.

\begin{proposition}\label{prop:thirteenpartial}
The Fermat curve $X_{13}$ has a nonconstant rational morphism to an
abelian fourfold $P_{13}$ such that
\[
 P_{13}(\Q)\simeq\Z/13\Z,\qquad
 2.632<L(P_{13}/\Q,1)<2.636.
\]
Its conductor is $13^8$, its archimedean factor is
$(2\pi)^{-4s}\Gamma(s)^4$, and its functional-equation sign is $+1$.
\end{proposition}
\begin{proof}
We give the quotient, rational-rank argument, Hecke decomposition,
and central-value certificate explicitly.

\paragraph{The two curve quotients.}
There is a degree-thirteen map
\[
 X_{13}\longrightarrow\mathcal K,
 \qquad [X:Y:Z]\longmapsto[X^4Z:Y^4X:Z^4Y],
 \quad\mathcal K:\ A^4B+B^4C+C^4A=0.
\]
Substitution multiplies the Fermat equation by $(XYZ)^4$.
On the torus, the exponent matrix has determinant thirteen, and its
kernel is generated by $[X:Y:Z]\mapsto[\zeta X:\zeta^4Y:Z]$.
The invariant functions
$x=-X^{13}/Z^{13}$ and $y=XY^3/Z^4$ satisfy
$y^{13}=x(1-x)^3$. Thus $\mathcal K$ is the smooth proper cyclic
quotient $C_{1,3,9}$, of genus six.

The rational automorphism $c[A:B:C]=[C:A:B]$ has order three.
Its two fixed points are $[1:\zeta_3:\zeta_3^2]$ and its conjugate;
$[1:1:1]$ is not on $\mathcal K$. Riemann--Hurwitz gives genus two
for $B=\mathcal K/\langle c\rangle$.
For an explicit equation put
\[
 S=A+B+C,\quad T=AB+BC+CA,\quad U=ABC,\quad
 D=(A-B)(B-C)(C-A).
\]
Direct expansion gives
\[
 2(A^4B+B^4C+C^4A)
 =S^3T-3ST^2-S^2U+5TU-D(S^2-T).
\]
With $t=T/S^2$, $u=U/S^3$ and $d=D/S^3$, the equations are
\[
 d(1-t)=t-3t^2+(5t-1)u,\qquad
 d^2=t^2-4t^3-4u-27u^2+18tu.
\]
Eliminating $d$ gives
\[
 (13t^2-16t+7)u^2+(-12t^3+14t^2-7t+1)u+t^5=0.
\]
Its discriminant is $(1-t)^2f(t)$, where
\[
 B:\ v^2=f(t)=1-12t+52t^2-104t^3+104t^4-52t^5.
\]
The quotient coordinate is
\[
 v=\frac{2(13t^2-16t+7)u-12t^3+14t^2-7t+1}{1-t}.
\]
These invariants separate cyclic orbits generically: $S,T,U$
determine the unordered triple and $D$ distinguishes its two
cyclic orientations. Consequently this equation gives the quotient
function field, not just a map to a genus-two curve.
The discriminant of $f$ is $2^8 13^3$.

\paragraph{The finite rational target.}
Let $J=\Jac(\mathcal K)$ and
$P_{13}=\ker(\operatorname{Nm}:J\to\Jac(B))^0$.
Pullback followed by norm is multiplication by three, so
$J\sim_\Q\Jac(B)\times P_{13}$ and $\dim P_{13}=4$.
The exact counts
\[
 \begin{array}{c|ccc}
 q&\#B(\F_q)&\#B(\F_{q^2})&\#\Jac(B)(\F_q)\\\hline
 3&7&11&27\\
 5&6&26&26
 \end{array}
\]
use $\#\Jac(B)(\F_q)=(\#B(\F_q)^2+\#B(\F_{q^2}))/2-q$.
Prime-to-residue-characteristic torsion injectivity at both good
primes shows $\Jac(B)(\Q)_{\rm tors}=0$. The distinct rational
points $\infty$ and $(0,1)$ therefore give a non-torsion divisor
class. On the other hand,
\citet[\S4, table on p.~219]{GrossRohrlich1978} proves
$\operatorname{rank}J(\Q)=1$. It follows that
$\operatorname{rank}\Jac(B)(\Q)=1$ and
$\operatorname{rank}P_{13}(\Q)=0$.

Their Theorem~1.1 gives $J(\Q)_{\rm tors}=\Z/13\Z$.
This torsion is contained in $P_{13}$: its norm is zero because
$\Jac(B)(\Q)$ has no torsion, and
$3x=(1-c)(2+c)x$ whenever $(1+c+c^2)x=0$.
Since multiplication by three is invertible on a group of order
thirteen, $x$ belongs to the image of $1-c$, namely $P_{13}$.
Thus $P_{13}(\Q)=\Z/13\Z$.
Composing an Abel--Jacobi map of $\mathcal K$ with $1-c$ gives a
nonconstant map to $P_{13}$: otherwise $1-c$ would vanish on the
Jacobian generated by the curve, contrary to $\dim P_{13}=4$.

\paragraph{The Hecke constituents.}
Put $L=\Q(\zeta_{13})$, $H=\langle3\rangle\subset(\Z/13)^\times$,
and $K_0=L^H$, a quartic CM field. The automorphism
\[
 \tau[A:B:C]=[\zeta_{13}^4 A:\zeta_{13}B:C]
\]
satisfies $c\tau c^{-1}=\tau^3$ and gives the usual faithful
$L$-action on $J$. Over an algebraic closure of a coefficient field,
the twelve nontrivial $\tau$-eigenspaces of $V_\ell J$ are lines.
The automorphism $c$ groups them into four triples; the invariant
part $V_\ell\Jac(B)$ has one dimension in each triple.
The central field $K_0$ preserves this invariant part and acts
faithfully on it. Its endomorphisms are defined over $K_0$, because
$G_{K_0}$ acts on $L$ through $H$, which fixes $K_0$.
Thus $\Jac(B)$ has CM by $K_0$, defined over $K_0$.

Choose a cubic Dirichlet character $\chi$ modulo thirteen with
$\chi(2)=\zeta_3$; since $2^4=3$ modulo thirteen, its restriction
to $H$ sends $3$ to $\zeta_3$.
For $g\in G_{K_0}$ whose action on $L$ is $\zeta_{13}\mapsto
\zeta_{13}^{3^e}$, the operator $gc^{-e}$ preserves every
$\tau$-eigenline and commutes with $c$. It therefore acts by the
same scalar on the three lines in each orbit. Passing to the
$c$-eigenspaces shows that their characters differ by $\chi^j$,
for $j=0,1,2$.
To verify this comparison for the full Galois group, put
$D=\{1,5,8,12\}$, so $(\Z/13)^\times=H\times D$.
Choose a vector $e_d$ in the $\tau$-eigenline with eigenvalue
$\zeta_{13}^d$ for each $d\in D$; the vectors $c^ke_d$
then lie in the lines indexed by $9^kd$.
If $g(\zeta_{13})=\zeta_{13}^a$ and
$a^{-1}=9^{k(a)}d(a)$, then
$g(e_d)=u(g,d)c^{k(a)}e_{d(a)d}$ for a scalar $u(g,d)$.
Because $g$ commutes with $c$, the basis
$\sum_{j=0}^2\lambda^{-j}c^je_d$ of the $c$-eigenspace
with eigenvalue $\lambda$ gives its matrix as the matrix on the
invariant space multiplied by $\lambda^{k(a)}$.
For $\lambda=\zeta_3$, this multiplier is $\chi(a)$.
Thus, after extending coefficients,
\begin{equation}\tag{5.8}\label{eq:thirteen-twists}
 V_\ell J\simeq V_\ell\Jac(B)\oplus
 (V_\ell\Jac(B)\otimes\chi)\oplus
 (V_\ell\Jac(B)\otimes\chi^2),
\end{equation}
and $V_\ell P_{13}$ is the sum of the last two terms.

The CM theorem identifies the one-dimensional $G_{K_0}$ characters
with algebraic Hecke characters; see
\citet[Theorems~3.5 and 3.7, pp.~10--11]{ConradCM}.
Induction from one of the four eigenspaces gives
$V_\ell\Jac(B)$ over $\Q$. The published Hecke factorization
and completed functional equation also follow directly from
\citet[Theorem~4 and Remark~2, pp.~187--188]{Milne1972}:
the CM field $K_0$ is Galois-stable in the endomorphism algebra,
its field of definition is $K_0$, and its four embeddings form
one Galois orbit. Therefore each twist in
\eqref{eq:thirteen-twists} has a degree-four Hecke $L$-function,
with two factors $\Gamma(s)$. The two twists have complex-conjugate
Dirichlet coefficients and dual functional equations. Their root
numbers are complex conjugates of modulus one, so their product
has root number $+1$. In particular $L(P_{13},s)$ is entire.

\paragraph{Conductor and local factors.}
The displayed genus-two model has good reduction away from two
and thirteen. At two use the integral model
\[
 w^2+w=-3t+13t^2-26t^3+26t^4-13t^5,
 \qquad v=2w+1,
\]
whose reduction has the smooth odd-degree Artin--Schreier form.
At thirteen, take $\pi^8=13$ and set $t=\pi^{-2}u$,
$v=\pi^{-1}y$. The equation becomes
\[
 y^2=\pi^2-12u+4\pi^6u^2-8\pi^4u^3+8\pi^2u^4-4u^5.
\]
Its smooth special fiber is $y^2=u-4u^5$. Thus the reduction is
potentially good over a tame extension. A tame generator
$\pi\mapsto\zeta_8\pi$ acts on the special fiber by
$(u,y)\mapsto(\zeta_8^2u,\zeta_8y)$.
On its differentials $du/y,u\,du/y$ the eigenvalues are
$\zeta_8,\zeta_8^3$; on first cohomology the four eigenvalues are
the primitive eighth roots. This can equally be checked on the
same smooth equation in characteristic zero and specialized.
There are no inertia invariants, so the conductor exponent of
$\Jac(B)$ is four. Twisting by $\chi$ or $\chi^2$ multiplies
these eigenvalues by third roots; none becomes one. Each twist
still has conductor exponent four and no wild contribution.
Consequently $P_{13}$ has conductor $13^8$ and Euler factor one
at thirteen. With $B_0=13^4/(2\pi)^4$ its completed function is
\[
 \Lambda_P(s)=B_0^s\Gamma(s)^4L(P_{13},s),\qquad
 \Lambda_P(s)=\Lambda_P(2-s).
\]

The program computes Euler polynomials of $J$ from the Jacobi sums
for $(1,3,9)$, and those of $\Jac(B)$ by exact point counts;
polynomial division gives the factors of $P_{13}$.
For example, the factors at two and three are
\[
 \begin{split}
 P_2(T)&=1-4T^4+16T^8,\\
 P_3(T)&=1-3T+4T^2+3T^3-11T^4+9T^5
                       +36T^6-81T^7+81T^8.
 \end{split}
\]
For $J$, the degree-three extension counts at three are
$\#\mathcal K(\F_{27})=55$ and
$\#\mathcal K(\F_{729})=913$; they give
$1+9T^3+71T^6+243T^9+729T^{12}$.
At primes of even residue degree modulo thirteen, the standard
Jacobi identity of \citet[Lemma~1.1]{GrossRohrlich1978} supplies
the factor of $J$. At split primes the coefficients are integer
Newton polynomials in the thirteen Jacobi residue counts.
Only the exact initial coefficients required by the cutoff are
used or reported.

\paragraph{A central-value certificate with a universal tail.}
Write $L(P_{13},s)=\sum a_nn^{-s}$. The same Mellin argument gives
\[
 L(P_{13},1)=2\sum_{n\ge1}\frac{a_n}{n}Q_4(n/B_0),\qquad
 Q_4(x)\le\Gamma(k+1)^4x^{-k}.
\]
There are eight reciprocal roots of absolute value $\sqrt p$ at
every good prime. Hence $|a_n|\le d_8(n)\sqrt n$, where
$\sum d_8(n)n^{-\sigma}=\zeta(\sigma)^8$.
For $\sigma>1$ and $k+1/2>\sigma$, the omitted Dirichlet tail
after $N$ is therefore at most
\begin{equation}\tag{5.9}\label{eq:thirteen-tail}
 2\Gamma(k+1)^4B_0^kN^{\sigma-k-1/2}\zeta(\sigma)^8.
\end{equation}
This bound uses only the degree-eight Weil estimate.

The finite sum is evaluated on $z=2+it$, $0\le t\le20$, as
\[
 \frac2\pi\operatorname{Re}\int_0^{20}
 \frac{\Gamma(3+it)^4}{2+it}B_0^{2+it}
 \sum_{n\le N}\frac{a_n}{n^{3+it}}\,dt.
\]
Seventeen gamma recurrences bound its vertical tail by
\[
 \frac2\pi B_0^2\left(\sum_{n\le N}\frac{|a_n|}{n^3}\right)
 \frac{\Gamma(20)^4}{68\cdot20^{68}}.
\]
For $N=100000$, Arb arithmetic at $144$ bits gives the finite
integral $2.633619912614\pm9.68\cdot10^{-14}$ and vertical tail
less than $2.612\cdot10^{-20}$.
Equation~\eqref{eq:thirteen-tail}, with $k=8$ and $\sigma=3/2$,
is less than $0.001459$. These enclosures prove the asserted
rational interval and complete the proof.
\end{proof}

The finite rational group and central nonvanishing just proved do
not by themselves identify the divisible part of
$\Sha(\Q,P_{13})$. In the quartic-CM description the two constituents
over $\Q(\sqrt{13})$ are nontrivial cubic twists of a CM Hilbert
form. Their central characters are cubic, so the trivial-central-
character version of Nekov\'a\v r's Theorem~B cannot be applied
directly. Section~\ref{subsec:lm-lz-rational} supplies the needed
arithmetic input by a different Euler-system argument: it proves
$T_\ell\Sha(\Q,P_{13})=0$ for all but finitely many $\ell$ splitting
completely in one fixed finite Galois extension. This is sufficient
for Theorem~\ref{thm:split-prime-descent}. The map constructed above
generates $P_{13}$ by differences, so its image satisfies the
geometric hypothesis of that theorem.

Passing to the real cubic field cut out by $\chi$ does not supply
this hypothesis through Zhang's rank-one theorem. To check this
precisely, restrict the four primitive eighth-root inertia characters
computed above to $F_0=\Q(\sqrt{13})$: the characters become
$i,-i,i,-i$. Real multiplication and the polarization give
cyclotomic determinant on each two-dimensional Hilbert constituent,
so its two distinct inertia characters are $i$ and $-i$.
Frobenius at thirteen fixes both, since $13\equiv1\pmod4$.
Thus each local Weil representation is a sum of two characters,
with zero monodromy, and the local automorphic representation is
principal series of conductor exponent two. Restriction to the
further cubic field preserves this property, while the totally
real degree becomes six. All other finite places have good reduction.
Consequently even a simple-zero certificate for the untwisted form
would not satisfy the finite-place hypothesis of
\citet[Theorem~A and its preceding setup, p.~28]{Zhang2001},
which in even degree requires a conductor exponent one, or the
nonprincipal-place alternative of
\citet[Theorem~4.3.2, p.~588]{Zhang2010}.
Nekov\'a\v r's Theorems~B and~B$'$ require central nonvanishing;
they are not rank-one statements.

The other full cyclic-quotient factors do not give a verified
rank-zero substitute. At exponent thirteen their three types are
$(1,1,11)$, $(1,2,10)$ and $(1,3,9)$, with multiplicities three,
six and two. All three have sign $-1$ in
\citet[table on p.~219]{GrossRohrlich1978}; the first two have
Hecke-character conductor $(1-\zeta_{13})^2$.
Their Theorem~2.1, pp.~208--209, and the same table prove rank one
for the last two types, but give only rank at most one for the first.
The table's footnote invokes Birch--Swinnerton-Dyer to predict
equality in that remaining case. No unconditional positive-rank
claim for that hyperelliptic factor is inferred here.
